\section{Signal Aspect Decision Logic}
\label{sec:interlocking:signal-decision}

Once a route has been successfully validated and reserved, the system determines the appropriate signal aspect to display. This decision logic enforces the safe representation of signal colors, ensuring that the displayed aspect always reflects both the status of the authorized route and the occupancy of downstream track sections. A signal will only clear if the associated route passes all safety checks and the next block is confirmed as free; otherwise, it remains at danger. Conditional interlocking is also considered, where dependencies on adjacent signals, track circuits, or point machines may influence the final decision.

The resulting output is expressed through the standard aspects of red, yellow, or green. Red represents the default and inherently safe condition, indicating that movement is not permitted. Yellow serves as a cautionary aspect, warning the driver that the train may proceed but must be prepared to stop at the next signal, often because the following block is occupied. Green conveys full clearance, authorizing the train to proceed at the permitted speed with assurance that the track ahead is safe. These updates are continuously processed by the Qt backend and rendered in real time on the graphical operator interface, providing both operators and drivers with immediate, unambiguous feedback on network status.
